\icode{N22X-MM-SI-L3-261013-001}
\nceeID[6]{2022湖南}{4}
1932年，查德威克用未知射线轰击氢核，发现这种射线是由质量与质子大致相等的中性粒子（即中子）组成。如图，中子以速度\(v_0\)分别碰撞静止的氢核和氮核，碰撞后氢核和氮核的速度分别为\(v_1\)和\(v_2\)。设碰撞为弹性正碰，不考虑相对论效应，下列说法正确的是\paren[]

\begin{tikzpicture}[>=stealth,font=\scriptsize]
	\begin{scope}
		\draw[dashed] (0,0)--(3,0);
		\shade[ball color=gray!60] (0,0) circle(.15);
		\shade[ball color=gray] (3,0) circle(.15);
		\draw[->] (.15,0)--node[above]{\(v_0\)}(1,0);
		\draw[->] (3.15,0)--node[above]{\(v_1\)}(4,0);
		\node at(0,-.15)[below]{中子};
		\node at(3,-.15)[below]{氢核};
	\end{scope}
	\begin{scope}[yshift=-1.5cm]
		\draw[dashed] (0,0)--(3,0);
		\shade[ball color=gray!60] (0,0) circle(.15);
		\shade[ball color=gray] (3,0) circle(.15);
		\draw[->] (.15,0)--node[above]{\(v_0\)}(1,0);
		\draw[->] (3.15,0)--node[above]{\(v_2\)}(4,0);
		\node at(0,-.15)[below]{中子};
		\node at(3,-.15)[below]{氦核};
	\end{scope}
\end{tikzpicture}
\begin{choices}
	\item 碰撞后氮核的动量比氢核的小
	\item 碰撞后氮核的动能比氢核的小
	\item \(v_2\)大于\(v_1\)
	\item \(v_2\)大于\(v_0\)
\end{choices}
\begin{answer}
   B\qquad 提示：\(v_2 = \frac{2m_1}{m_1 + m_2}v_0\) \quad \(p_2 = m_2v_2 = \frac{2m_2}{m_1 + m_2}\cdot m_1v_0\)\quad \(E_{k2} = \frac{1}{2}m_2v_2^2 = \frac{m_1m_2}{(m_1 + m_2)^2}\cdot \frac{1}{2}m_1v_0^2 = \frac{1}{\frac{m_1}{m_2} + 2 + \frac{m_2}{m_1}}\cdot \frac{1}{2}m_1v_0^2\)
\end{answer}